% !TEX root = ../BeamerTemplate.tex
%%%%%%%%%%%%%%%%%%%%%%%%%%%%%%%%%%%%%%%%%%%%%%%%%%%%%%%%%%%%%%%%%%%%%%%%%%%%%%%%%%

\section{MIMO Communication Performance Evaluation}
\note{
  \begin{itemize}
    \item In this section, we will evaluate the MIMO communication performance of the lens-assisted and baseline systems.
    \item We will present the simulation configuration, the paired lens and baseline trajectories, and the evaluation metrics used to assess the system performance.
    \item We will also discuss the results of the paired lens and baseline trajectories, as well as the higher-order lens-assisted scaling.
  \end{itemize}
}

\subsection{Paired Lens and Baseline Trajectories}

\begin{frame}[t]{Two Controlled Protocols Extend the Component Study}
\begin{columns}[T,onlytextwidth]
  \column{0.58\textwidth}
  \centering
  \includegraphics[width=\linewidth,height=4.05cm,keepaspectratio]
    {Analisis_v4_ISO_TX_Lens_vs_WithoutLens_Linear_HalfCircular_3x3/figures_matlab/01_geometry_and_arrays_matlab.png}
  \ThesisFigureCaption{3-1}{Matched TX trajectories and the two complete RX designs.}

  \column{0.38\textwidth}
  \scriptsize
  \begin{block}{Protocol A: placement comparison}
    Focal-arc lens RX versus linear baseline, both $3\times3$, at identical
    linear and half-circular waypoints.
  \end{block}
  \begin{block}{Protocol B: order comparison}
    Complete lens-assisted $3\times3$, $5\times5$, and $9\times9$ systems on one
    linear trajectory.
  \end{block}
  {\tiny\color{red!70!black}\textbf{Boundary:} placement and imported RX pattern
  change together.\par}
\end{columns}
\SourceNote{Chapter 3: Simulation Configuration and Paired Trajectories}
\note{
  \begin{itemize}
    \item In this section we test Antenna-lens-assisted MIMO communication performance.
    \item We will present the simulation configuration, the paired lens and baseline trajectories, and the evaluation metrics used to assess the system performance.
    \item We will also discuss the results of the paired lens and baseline trajectories, as well as the higher-order lens-assisted scaling.
  \end{itemize}
}
\end{frame}


\begin{frame}[t]{System-Level Simulation Configuration}
\scriptsize
\begin{columns}[T,onlytextwidth]
  \column{0.53\textwidth}
  \renewcommand{\arraystretch}{1.08}
  \begin{tabularx}{\linewidth}{@{}>{\bfseries}p{0.43\linewidth}X@{}}
    \toprule
    Parameter & Value\\
    \midrule
    Room / carrier & $20\times20\times3$~m / 38~GHz\\
    Bandwidth & 400~MHz, 401 bins\\
    TX array & 3 isotropic elements, 10~dBm each\\
    Mobility & 1~m/s; 16 time samples/waypoint\\
    Ray samples & 100{,}000 per source\\
    Linear path & 21 waypoints; 19.00~m\\
    Half-circle & 37 waypoints; 9~m radius\\
    \bottomrule
  \end{tabularx}

  \column{0.43\textwidth}
  \begin{block}{RX branch populations}
    Linear: $21\times3=63$ waypoint--RX observations per design.\par\smallskip
    Half-circle: $37\times3=111$ observations per design.
  \end{block}
  \begin{block}{Pattern contrast}
    Baseline branches reuse the $0^\circ$ patch pattern; lens branches use three
    angle-specific patterns at $+45^\circ$, $0^\circ$, and $-45^\circ$.
  \end{block}
  \MethodNote{Matching indices are ordinal design pairs, not co-located
  antenna elements.}
\end{columns}
\SourceNote{Chapter 3: Common Simulation Configuration}
\note{
  \begin{itemize}
    \item This parameter table shows the simulation configuration for the lens-assisted and baseline systems.
  \end{itemize}
}
\end{frame}


\begin{frame}[t]{Three Metric Scopes Answer Different Questions}
\scriptsize
\begin{columns}[T,onlytextwidth]
  \column{0.32\textwidth}
  \begin{block}{1. Link budget}
    \centering
    \begin{equation*}
      \begin{aligned}
        P_{\rm RX}&=P_{\rm TX}\sum_i|H_i|^2,\\[-1pt]
        C_{\rm link}&=\mathbb E_f[\log_2(1+\gamma)]
      \end{aligned}
      \tag{3.1}
    \end{equation*}
    \vspace{-0.4em}
    \raggedright
    Three simultaneous TX contributions feed one active RX branch. This retains
    absolute path loss and pattern gain.
  \end{block}

  \column{0.32\textwidth}
  \begin{block}{2. Constructed MIMO}
    \centering
    $\mathbf H_{N\times N}\leftarrow$ sequential $N\times1$ runs\par\smallskip
    \raggedright
    Condition number, effective rank, and equal-power capacity at a fixed
    normalized SNR of 10~dB diagnose channel structure.
  \end{block}

  \column{0.32\textwidth}
  \begin{block}{3. Spatial decorrelation}
    \centering
    \begin{equation*}
      D_{ij}=1-|\rho_{ij}|
      \tag{3.4}
    \end{equation*}
    \vspace{-0.4em}
    \raggedright
    Center-TX responses are evaluated by global pooling, waypoint-local
    pooling, and block-median statistics.
  \end{block}
\end{columns}
\CautionNote{These metrics have different normalization and weighting. In
particular, $C_{\rm link}$ and fixed-10-dB constructed-MIMO capacity are not
numerically interchangeable.}
\SourceNote{Chapter 3: Evaluation Metrics and Aggregation}
\note{
  \begin{itemize}
    \item There are 3 metric scopes to answer different questions: link budget, constructed MIMO, and spatial decorrelation.
    \item For the constructed-MIMO scope, each sequential $N\times1$ channel
    realization forms one column of
    $\mathbf H(f)=[\mathbf h_1(f)\ \mathbf h_2(f)\ \cdots\ \mathbf h_N(f)]$,
    producing a constructed $N\times N$ channel matrix.
    \item From its singular values, the condition number and effective rank are
    $\kappa(f)=\sigma_1(f)/\sigma_3(f)$ and
    $r_{\mathrm{eff}}(f)=\exp[-\sum_{i=1}^{3}p_i(f)\ln p_i(f)]$, where
    $p_i(f)=\sigma_i^2(f)/\sum_{j=1}^{3}\sigma_j^2(f)$, as given in
    Equation~(3.2).
    \item The equal-power constructed-MIMO capacity is
    $C_{\mathrm{MIMO}}(\gamma)=\mathbb E_f[\sum_{i=1}^{3}
    \log_2(1+\gamma\sigma_i^2(f)/N_t)]$, with $N_t=3$, as given in
    Equation~(3.3).
    \item Emphasize that the matrix combines sequential simulations and is used
    to diagnose channel structure and diversity potential; it is not a directly
    measured simultaneous three-port channel.
  \end{itemize}
}
\end{frame}


\begin{frame}[t]{A Static Snapshot Shows Gain --- but Mixed MIMO Structure}
\begin{columns}[T,onlytextwidth]
  \column{0.59\textwidth}
  \centering
  \includegraphics[width=\linewidth,height=4.95cm,keepaspectratio]
    {Analisis_v4_ISO_TX_Lens_vs_WithoutLens_Linear_HalfCircular_3x3/figures_matlab/02_static_rx_scenario_comparison_matlab.png}
  \ThesisFigureCaption{3-2}{Static branch gain, TX correlation, and normalized capacity.}

  \column{0.37\textwidth}
  \scriptsize
  \begin{block}{Absolute branch gain}
    Lens branches span $-76.18$ to $-72.98$~dB; baseline spans only
    $-79.84$ to $-79.75$~dB.
  \end{block}
  \begin{block}{Constructed $3\times3$ snapshot}
    \renewcommand{\arraystretch}{1.12}
    \begin{tabular}{@{}lcc@{}}
      & Base & Lens\\\midrule
      $\kappa$ (dB) & \textbf{23.90} & 28.70\\
      $r_{\rm eff}$ & 1.42 & \textbf{1.60}\\
      $C_{\rm MIMO}$ & 6.81 & \textbf{7.24}\\
      RX corr. & 0.7927 & \textbf{0.5636}
    \end{tabular}
  \end{block}
  % \CautionNote{A better condition number alone does not prove stronger
  % multiplexing.}
\end{columns}
\SourceNote{Chapter 3: Static Snapshot Characterization}
\note{
  \begin{itemize}
    \item This snapshot shows the performance of the lens-assisted and baseline systems at a static point in time.
    \item The lens has better absolute branch gain, but the constructed MIMO metrics show mixed results.
    \item The lens improves rank, normalized capacity, and RX correlation in this snapshot, while its condition number becomes worse.
    \item This is a normalized single-position result. It does not imply that
          every lens branch retains a gain or capacity advantage as the TX moves
          along either trajectory.
  \end{itemize}
}
\end{frame}


\begin{frame}[t]{Equal-Weight Trajectory Aggregation Favors the Baseline}
\scriptsize
\centering
\renewcommand{\arraystretch}{1.18}
\begin{tabular}{@{}llrrrrrr@{}}
  \toprule
  Path & RX & $|H|_{\rm med}$ & $C_{\rm link}$ & $\kappa$ & $r_{\rm eff}$ & $C_{\rm MIMO}$ & $D_{\rm med}$\\
  \midrule
  Linear & Base & $-79.74$ & \textbf{5.41} & \textbf{25.01} & \textbf{1.202} & \textbf{6.12} & 0.0909\\
  Linear & Lens & $-86.75$ & 3.57 & 25.55 & 1.154 & 5.85 & \textbf{0.5308}\\
  Half-circle & Base & $-80.89$ & \textbf{5.08} & 25.58 & \textbf{1.348} & \textbf{6.71} & 0.0638\\
  Half-circle & Lens & $-87.29$ & 3.53 & \textbf{21.56} & 1.305 & 6.58 & \textbf{0.5518}\\
  \bottomrule
\end{tabular}

\vspace{0.35cm}
\begin{columns}[T,onlytextwidth]
  \column{0.48\textwidth}
  \MethodNote{\textbf{Aggregate rate and rank:} all three directional RX cases
  receive equal statistical weight, including persistently misaligned beams.}
  \column{0.48\textwidth}
  \MethodNote{\textbf{Local angular differentiation:} block-median
  decorrelation strongly favors the lens on both paths.}
\end{columns}
% \Takeaway{The lens improves selectivity, but equal weighting does not convert
% that selectivity into higher aggregate capacity.}
\SourceNote{Chapter 3: Aggregate Placement-Comparison Results; capacities in bit/s/Hz}
\note{
  \fontsize{7}{8}\selectfont
  \begin{itemize}
    \setlength{\itemsep}{0.3em}
    \setlength{\parskip}{0pt}
    \item This table compares the lens-assisted and baseline RX systems along
          the linear and half-circular trajectories.
    \item Equal-weight aggregation gives the same statistical weight to all
          three directional RX branches, including weak or misaligned lens
          branches; no beam selection is applied.
    \item Under this rule, the baseline achieves higher link capacity,
          effective rank, and MIMO capacity on both trajectories. This result is
          consistent with its more uniform branch response as the transmitter
          moves; it does not isolate antenna pattern from RX placement.
    \item The lens nevertheless provides much stronger median-block
          decorrelation, confirming useful local angular selectivity.
    \item Therefore, the lens is not universally worse. Its selectivity must be
          exploited through beam selection or adaptive combining to avoid
          persistently misaligned branches and potentially improve capacity.
    \item The takeaway is that the lens improves selectivity, but equal weighting does not convert that selectivity into higher aggregate capacity.
  \end{itemize}
}
\end{frame}


\begin{frame}[t]{Along Both Paths, Branch Alignment Drives a Wide Lens Spread}
\centering
\includegraphics[width=0.92\linewidth,height=4.55cm,keepaspectratio]
  {Analisis_v4_ISO_TX_Lens_vs_WithoutLens_Linear_HalfCircular_3x3/figures_matlab/04_channel_metrics_along_trajectory_matlab.png}
\ThesisFigureCaption{3-5}{Mean SNR and equivalent-SISO capacity; shaded lens
bands are RX-branch minima--maxima.}
\begin{columns}[T,onlytextwidth]
  \column{0.48\textwidth}
  \MethodNote{The baseline mean is higher at all 21 linear waypoints and at 28
  of 37 half-circle waypoints; aligned lens branches reverse the mean locally.}
  \column{0.48\textwidth}
  \CautionNote{All cases have 0\% outage at 1~bit/s/Hz; this threshold is
  saturated and cannot distinguish reliability.}
\end{columns}
\note{
  \fontsize{7}{8}\selectfont
  \begin{itemize}
    \setlength{\itemsep}{0.3em}
    \setlength{\parskip}{0pt}
    \item The solid and dashed lines are means across the three RX design
          cases. The shaded lens region is the minimum-to-maximum branch range,
          not a confidence interval.
    \item On the linear path, the baseline mean capacity is higher at all 21
          waypoints. On the half-circle, it is higher at 28 of 37 waypoints.
    \item The remaining half-circle positions show that a locally aligned lens
          branch can raise the three-branch mean above the baseline. The lens is
          therefore not weaker at every position, even though its aggregate
          median is lower.
    \item The much wider lens band is evidence of angular selectivity: one
          branch may align strongly while another remains weak at the same
          waypoint.
    \item Every configuration records 0\% outage at 1~bit/s/Hz. This threshold
          is saturated, so it cannot support a reliability comparison; a higher
          threshold or capacity percentiles would be more informative.
  \end{itemize}
}
\end{frame}


\begin{frame}[t]{The Constructed MIMO Channel Has a Geometry Bottleneck}
\centering
\includegraphics[width=0.91\linewidth,height=4.35cm,keepaspectratio]
  {Analisis_v4_ISO_TX_Lens_vs_WithoutLens_Linear_HalfCircular_3x3/figures_matlab/05_trajectory_mimo_metrics_matlab.png}
\ThesisFigureCaption{3-6}{Conditioning, effective rank, and fixed-10-dB
constructed-MIMO capacity.}
\begin{columns}[T,onlytextwidth]
  \column{0.48\textwidth}
  \MethodNote{At the 9.5~m linear midpoint, $\kappa$ reaches 65.69~dB
  (baseline) and 55.09~dB (lens).}
  \column{0.48\textwidth}
  \MethodNote{The half-circle avoids this isolated symmetry-related impairment
  and gives higher median rank and capacity for both designs.}
\end{columns}
\SourceNote{Chapter 3: Constructed MIMO Performance Along the Trajectories}
\note{
  \fontsize{7}{8}\selectfont
  \begin{itemize}
    \setlength{\itemsep}{0.3em}
    \setlength{\parskip}{0pt}
    \item This figure compares the constructed-MIMO conditioning, effective
          rank, and fixed-10-dB capacity along both trajectories.
    \item Under the evaluated equal-weight configuration, the lens generally
          has lower effective rank and MIMO capacity than the baseline. Its
          angular selectivity does not become a multiplexing gain when weak or
          misaligned branches remain included.
    \item The condition-number result is mixed: the lens is 0.55~dB worse in
          median on the linear path but 4.02~dB better on the half-circle.
    \item At the 9.5~m linear midpoint, geometric symmetry creates an isolated
          bottleneck. The condition number reaches 65.69~dB for the baseline
          and 55.09~dB for the lens. The lens reduces the severity, but both
          channels remain extremely ill-conditioned.
    \item The half-circle avoids this symmetry-related impairment and gives
          higher median rank and capacity for both designs, showing that
          deployment geometry remains a dominant factor.
    \item Therefore, the lens is not worse for every metric, but it
          underperforms the baseline in the main multiplexing metrics under
          this operating rule. Beam selection or adaptive combining is needed
          to exploit its directional response.
  \end{itemize}
}
\end{frame}


\begin{frame}[t]{Lens Decorrelation Gain Depends on Aggregation Scope}
\begin{columns}[T,onlytextwidth]
  \column{0.61\textwidth}
  \centering
  \includegraphics[width=\linewidth,height=4.35cm,keepaspectratio]
    {Analisis_v4_ISO_TX_Lens_vs_WithoutLens_Linear_HalfCircular_3x3/figures_matlab/06_spatial_decorrelation_comparison_matlab.png}
  \ThesisFigureCaption{3-7}{Waypoint-local decorrelation of the center-TX response.}

  \column{0.35\textwidth}
  \scriptsize
  \begin{block}{Global pooled}
    Lens rises only slightly: 0.9784 to 0.9807 (linear) and 0.9744 to 0.9872
    (half-circle).
  \end{block}
  \begin{block}{Typical local block}
    Lens $D_{\rm med}$: 0.5308 / 0.5518.\par
    Baseline: 0.0909 / 0.0638.
  \end{block}
  {\tiny\color{red!70!black}\textbf{Not equivalent to} the singular-value
  structure of the constructed $3\times3$ matrix.\par}
\end{columns}
% \Takeaway{Both estimators favor the lens; aggregation scope determines whether
% the measured increase is slight or substantial.}
\SourceNote{Chapter 3: Spatial Decorrelation}
\note{
  \fontsize{7}{8}\selectfont
  \begin{itemize}
    \setlength{\itemsep}{0.3em}
    \setlength{\parskip}{0pt}
    \item This slide shows that the decorrelation result depends on the
          aggregation scope and on where normalization is applied. A higher
          $D=1-|\rho|$ means less-correlated RX responses.
    \item Global pooling concatenates all waypoint, time, and frequency samples
          before computing correlation. Under this estimator the lens is only
          slightly higher on both trajectories.
    \item The block-median estimator evaluates correlation locally before
          taking a typical value. It strongly favors the lens: $D_{\rm med}$
          increases from 0.0909 to 0.5308 on the linear path and from 0.0638 to
          0.5518 on the half-circle.
    \item The effect sizes differ because global pooling measures the coherence
          of the complete concatenated trajectory, whereas the local estimator
          shows that a typical lens block has more differentiated
          angle-dependent responses.
    \item Thus, aggregation changes the magnitude of the observed lens benefit,
          not its direction: both global-pooled and block-median decorrelation
          are higher with the lens in this $3\times3$ comparison.
    \item This calculation uses only the center-TX response across three RX
          branches. It is not equivalent to the singular values, effective
          rank, or multiplexing performance of the constructed $3\times3$
          MIMO matrix.
    \item The appropriate conclusion is that the lens provides stronger local
          angular selectivity, not that it automatically provides higher MIMO
          capacity.
    \item Takeaway: both estimators favor the lens; aggregation scope determines whether the measured increase is slight or substantial.
  \end{itemize}
}
\end{frame}


\begin{frame}[t]{System-Level Interpretation: Select, Combine, or Pay the Penalty}
\centering
\resizebox{0.92\linewidth}{!}{%
\begin{tikzpicture}[>=latex,font=\footnotesize,node distance=0.70cm]
  \node[draw=Blue,fill=SoftBlue,rounded corners,minimum width=3.3cm,
        minimum height=1.0cm,align=center] (lens) {lens produces\\angle-specific branches};
  \node[draw=Teal,fill=SoftTeal,rounded corners,minimum width=3.3cm,
        minimum height=1.0cm,align=center,right=1.2cm of lens] (spread) {large within-waypoint\\capacity spread};
  \node[draw=Gold!85!black,fill=SoftGold,rounded corners,minimum width=3.3cm,
        minimum height=1.0cm,align=center,right=1.2cm of spread] (decision) {receiver operating rule};
  \node[draw=red!70!black,rounded corners,minimum width=4.0cm,
        minimum height=1.0cm,align=center,below left=1.25cm and -0.1cm of decision] (equal) {equal weighting\\misaligned branches reduce aggregate rate};
  \node[draw=Teal,rounded corners,minimum width=4.0cm,
        minimum height=1.0cm,align=center,below right=1.25cm and -0.1cm of decision] (adapt) {beam selection / combining\\can exploit the strongest local branch};
  \draw[->,thick] (lens)--(spread);
  \draw[->,thick] (spread)--(decision);
  \draw[->,thick] (decision)--(equal);
  \draw[->,thick] (decision)--(adapt);
\end{tikzpicture}}

\vspace{0.45cm}
\begin{columns}[T,onlytextwidth]
  \column{0.48\textwidth}
  \CautionNote{Measured result in this thesis: equal-weight lens aggregation
  yields lower median link and constructed-MIMO capacity.}
  \column{0.48\textwidth}
  \MethodNote{Design implication: the lens is an angularly selective front end,
  not a uniform-gain replacement for a coplanar array.}
\end{columns}
\SourceNote{Chapter 3: Discussion and Placement-Comparison Summary}
\note{
  \fontsize{7}{8}\selectfont
  \begin{itemize}
    \setlength{\itemsep}{0.3em}
    \setlength{\parskip}{0pt}
    \item The system-level conclusion is that the lens behaves as an
          angle-selective front end, rather than as a uniform-gain replacement
          for a conventional coplanar array.
    \item Its directional branches produce a wide capacity spread at each
          waypoint: an aligned branch can be strong, while misaligned branches
          can remain weak.
    \item The receiver operating rule therefore determines whether this
          selectivity becomes useful. Giving every branch equal weight retains
          the weak branches and reduces the aggregate link and constructed-MIMO
          capacity in the results of this thesis.
    \item Beam selection or adaptive combining could instead emphasize the
          strongest locally aligned branch and avoid persistently misaligned
          responses.
    \item However, this adaptive strategy is a design implication and future
          research direction; it was not directly evaluated in the present
          equal-weight simulations.
    \item The appropriate conclusion is not that the lens is universally
          better or worse, but that its benefit is conditional on alignment
          and on how the receiver processes the angle-specific branches.
  \end{itemize}
}
\end{frame}


\subsection{Higher-Order Lens-Assisted Scaling}

\begin{frame}[t]{Higher-Order Study Changes Power, Aperture, and Angular Grid}
\begin{columns}[T,onlytextwidth]
  \column{0.61\textwidth}
  \centering
  \includegraphics[width=\linewidth,height=4.50cm,keepaspectratio]
    {Analisis_v4_ISO_TX_Lens_Order_3x3_5x5_9x9_Linear/figures_matlab/01_geometry_order_comparison_matlab.png}
  \ThesisFigureCaption{3-9}{Common linear path with order-dependent TX and RX geometry.}

  \column{0.35\textwidth}
  \scriptsize
  \setlength{\tabcolsep}{2pt}
  \renewcommand{\arraystretch}{1.10}
  \begin{tabular}{@{}lrrr@{}}
    \toprule
    $N$ & Aperture & RX step & $P_{\rm TX,tot}$\\
    \midrule
    3 & 0.90~m & $45^\circ$ & 14.77~dBm\\
    5 & 1.80~m & $30^\circ$ & 16.99~dBm\\
    9 & 3.60~m & $15^\circ$ & 19.54~dBm\\
    \bottomrule
  \end{tabular}

  \vspace{0.30cm}
  \begin{block}{Normalization}
    \[
      \eta_r=\frac{r_{\rm eff}}{N},\qquad
      \eta_C=\frac{C_{\rm MIMO}}{N}
    \]
  \end{block}
  \CautionNote{This is complete-design scaling, not a causal test of order $N$
  alone.}
\end{columns}
\SourceNote{Chapter 3: Higher-Order Simulation Configuration}
\end{frame}


\begin{frame}[t]{Higher Order Raises Capacity but Lowers Efficiency}
\scriptsize
\begin{columns}[T,onlytextwidth]
  \column{0.47\textwidth}
  \renewcommand{\arraystretch}{1.18}
  \begin{tabularx}{\linewidth}{@{}Xrrr@{}}
    \toprule
    Metric & $3\times3$ & $5\times5$ & $9\times9$\\
    \midrule
    $C_{\rm link}$ & 3.57 & 4.13 & \textbf{5.19}\\
    $\kappa$ (dB) & \textbf{25.55} & 34.24 & 42.88\\
    $r_{\rm eff}$ & 1.154 & 1.332 & \textbf{2.043}\\
    $\eta_r$ & \textbf{0.385} & 0.266 & 0.227\\
    $C_{\rm MIMO}$ & 5.851 & 8.254 & \textbf{14.562}\\
    $\eta_C$ & \textbf{1.950} & 1.651 & 1.618\\
    $D_{\rm global}$ & \textbf{0.9807} & 0.9276 & 0.8570\\
    $D_{\rm med}$ & \textbf{0.5308} & 0.5166 & 0.4876\\
    \bottomrule
  \end{tabularx}
  \SlideCaption{Trajectory medians; capacities in bit/s/Hz}

  \column{0.49\textwidth}
  \centering
  \includegraphics[width=\linewidth,height=4.85cm,keepaspectratio]
    {Analisis_v4_ISO_TX_Lens_Order_3x3_5x5_9x9_Linear/figures_matlab/08_aggregate_order_comparison_matlab.png}
  \ThesisFigureCaption{3-12}{Absolute, normalized, and decorrelation metrics.}
\end{columns}
% \Takeaway{$3\times3\rightarrow9\times9$: fixed-10-dB capacity rises 148.9\%,
% but rank fraction falls 41.0\% and capacity per dimension falls 17.0\%.}
\SourceNote{Chapter 3: Aggregate Higher-Order Results}
\note{
  \fontsize{7}{8}\selectfont
  \begin{itemize}
    \setlength{\itemsep}{0.3em}
    \setlength{\parskip}{0pt}
    \item This slide separates the benefit of adding dimensions from the
          efficiency with which those dimensions are used.
    \item From $3\times3$ to $9\times9$, link capacity rises from 3.57 to
          5.19~bit/s/Hz, and fixed-10-dB constructed-MIMO capacity rises by
          148.9\%, from 5.851 to 14.562~bit/s/Hz.
    \item However, conditioning worsens as $\kappa$ increases from 25.55 to
          42.88~dB. Although absolute effective rank increases, the rank
          fraction $\eta_r$ falls by 41.0\%, showing that a smaller fraction of
          the available dimensions remains effective.
    \item Capacity per dimension $\eta_C$ also falls by 17.0\%, and both
          decorrelation estimators decrease. The additional dimensions are
          therefore increasingly correlated and less efficiently utilized.
    \item The conclusion is that higher order increases total capacity, but
          with diminishing spatial efficiency. This is complete-design scaling,
          not an order-only effect, because transmit power, aperture, and the
          angular grid also change with $N$.
    \item Takeaway : {$3\times3\rightarrow9\times9$: fixed-10-dB capacity rises 148.9\%,
but rank fraction falls 41.0\% and capacity per dimension falls 17.0\%.}
  \end{itemize}
}
\end{frame}


\begin{frame}[t]{More Dimensions Do Not Remove Geometry-Limited Modes}
\begin{columns}[T,onlytextwidth]
  \column{0.62\textwidth}
  \centering
  \includegraphics[width=\linewidth,height=4.05cm,keepaspectratio]
    {Analisis_v4_ISO_TX_Lens_Order_3x3_5x5_9x9_Linear/figures_matlab/05_trajectory_mimo_scaling_matlab.png}
  \ThesisFigureCaption{3-14}{Order-dependent conditioning, rank efficiency, and capacity.}

  \column{0.34\textwidth}
  \scriptsize
  \begin{block}{Shared midpoint bottleneck}
    At 9.5~m, $r_{\rm eff}\approx1$ for every order; $\kappa$ worsens from
    55.09 to 71.89~dB.
  \end{block}
  \begin{block}{Decorrelation trend}
    Global and block-median decorrelation decrease as the RX pair population
    grows and angular spacing tightens.
  \end{block}
  {\tiny\color{red!70!black}\textbf{Confounds:} +4.77~dB total TX power,
  fourfold aperture, and only one common RX angle.\par}
\end{columns}

\SourceNote{Chapter 3: Constructed MIMO Scaling and Spatial Decorrelation}
\note{
  \fontsize{7}{8}\selectfont
  \begin{itemize}
    \setlength{\itemsep}{0.3em}
    \setlength{\parskip}{0pt}
    \item This figure shows that adding more antenna dimensions does not remove
          a geometry-induced channel bottleneck.
    \item At the 9.5~m midpoint, the effective rank collapses to approximately
          one for every order. The condition number also worsens from 55.09~dB
          for $3\times3$ to 71.89~dB for $9\times9$.
    \item Thus, when the propagation geometry makes the channel responses
          nearly dependent, a larger array adds nominal dimensions but does not
          create an equal number of usable spatial modes.
    \item Global and block-median decorrelation decrease with order, consistent
          with greater similarity among the growing RX-pair population and the
          tighter angular spacing.
    \item This trend cannot be attributed to $N$ alone: total TX power rises by
          4.77~dB, the aperture changes by a factor of four, and the compared
          angular grids share only one common RX angle.
    \item The safe conclusion is that higher order raises absolute capacity,
          but geometry-limited modes remain and spatial efficiency decreases.
    \item Takeaway : Higher order adds capacity, but its dimensions are less balanced and
less efficiently utilized.
  \end{itemize}
}
\end{frame}
